\section{Per-POI Noise-Floor Tables (Phase-0 Detail)}
\label{app:noise-floor-tables}

This appendix tabulates the per-POI noise-floor distributions from Exp~2 (Section~\ref{sec:experiments-baseline}) on the 10-POI probe set sorted by $|\nfmu|$. Each row reports $\nfmu$, $\nflo$, $\nfhi$ under $N{=}200$ random size-matched $\approx 80$~Da phospho-mass perturbations, plus the observed $\dptraw$ for the corresponding probe-set positive at the experimentally annotated phospho site. Clearance is evaluated against the pre-registered $\nftau{=}1.5$ threshold (Eq.~\ref{eq:dpt-cal}). $\simmark$

\begin{table}[h]
\centering
\caption{Phase-0 per-POI noise-floor table on the 10-POI probe set. The lower-$|\nfmu|$ POIs clear the threshold; the top three by $|\nfmu|$ miss because their per-POI noise floor is wider than the phospho-driven gap. The MD-relaxed fallback recovers 2 of the 3 (\S\ref{sec:experiments-baseline}). $\simmark$}
\label{tab:appa-nf}
\begin{tabular}{lccccc}
\toprule
POI (probe-set rank) & $|\nfmu|$ & $[\nflo, \nfhi]$ width & $\dptraw$ at PTM site & $\nftau|\nfmu|$ & Clearance \\
\midrule
POI-1 (BCL-XL-phospho-S62) & $0.024$ & $0.18$ & $+0.140$ & $0.036$ & $\checkmark$\\
POI-2 (BCL-2-phospho-T56) & $0.026$ & $0.16$ & $+0.120$ & $0.039$ & $\checkmark$\\
POI-3 (PCNA-phospho-Y211) & $0.027$ & $0.16$ & $+0.110$ & $0.040$ & $\checkmark$\\
POI-4 (FOXO3a-phospho-S253) & $0.030$ & $0.14$ & $-0.100$ & $0.045$ & $\checkmark$\\
POI-5 (CDK1-phospho-T161) & $0.031$ & $0.14$ & $+0.090$ & $0.046$ & $\checkmark$\\
POI-6 (BAD-phospho-S136) & $0.034$ & $0.12$ & $+0.100$ & $0.050$ & $\checkmark$\\
POI-7 (HDAC4-phospho-S246) & $0.037$ & $0.12$ & $-0.090$ & $0.055$ & $\checkmark$\\
POI-8 (cMyc-phospho-T58) & $0.052$ & $0.10$ & $+0.060$ & $0.078$ & $\times$ \\
POI-9 (BAX-phospho-S184) & $0.058$ & $0.10$ & $+0.050$ & $0.087$ & $\times$ \\
POI-10 (MCL-1-phospho-T163) & $0.061$ & $0.08$ & $+0.040$ & $0.092$ & $\times$ \\
\bottomrule
\multicolumn{6}{l}{\footnotesize POI-8, POI-9, POI-10 have low Boltz-2 pLDDT $\leq 67$ in the modification region; MD fallback recovers POI-8 and POI-9.}\\
\end{tabular}
\end{table}

\paragraph{Reading the table.} The clearance criterion in Eq.~\ref{eq:dpt-cal} requires $|\dptraw| > \nftau|\nfmu|$. The first seven probe-set POIs satisfy this comfortably; the last three (POI-8 through POI-10) sit inside their per-POI noise floor and are routed to the zero branch. These three POIs share a common failure mode: the Boltz-2 single-state $\poiptm$ prediction has low local pLDDT at the PTM site, which widens the noise floor by injecting structural uncertainty into the perturbation distribution. The MD-relaxed fallback (Stage B$'$ in Figure~\ref{fig:pipeline}) reduces local structural uncertainty at the cost of $\approx 8$~hours of MD wall-clock per tuple; on POI-8 and POI-9 the fallback brings the $|\dptraw|$ above the cleared threshold by tightening the WT structure.

\paragraph{Mass-matched calibration per PTM type.} The same protocol is run with PTM-mass-matched perturbations for the cross-track experiments (\S\ref{sec:experiments-scope}): $\approx 8.5$~kDa for mono-ubiquitylation, $14$--$42$~Da for methylation, $\approx 42$~Da for acetylation. The methylation noise-floor magnitudes are reported in Appendix~\ref{app:cross-track-detail}.
